\section{Limitations}\label{sec:limitations}

Five limitations define the scope of the present evidence. First, the
experiments labelled as GRU were not executed with a true recurrent
backend; the audit identified fallback multilayer-perceptron execution
on flattened sequences. These runs are therefore useful as sequentially
labelled stress tests, but they cannot support claims about recurrent
neural networks or long short-term memory learning in rainfall--runoff
modelling \citep{kratzert2018rainfall}. Second, the stochastic ensemble
is not yet a calibrated predictive-uncertainty product. Its empirical
coverage is far below the nominal 95\% level, so its current value is as
a source of distributional descriptors for ML correction rather than as
an operational uncertainty band. Coverage recalibration, bias-aware
post-processing and formal reliability assessment would be required
before operational use \citep{thiboult2016accounting}.

Third, the low-flow problem remains only partially resolved. The linear
baseflow reservoir improves the global deterministic baseline, but the
regime diagnostics show that it does not remove low-flow RMSE and bias.
Future calibration should therefore include objective functions that
explicitly weight low flows, such as inverse-flow criteria, rather than
relying only on global NSE, KGE or RMSE \citep{pushpalatha2012review}.
Fourth, the physical component is lumped and was evaluated in one
strongly seasonal high-Andean basin. The conclusions should therefore be
read as evidence for this hydroclimatic and modelling context, not as a
claim of universal superiority over distributed hydrological models or
over ML-only approaches in all basins. Fifth, the audit status is
reported as WARN rather than fully PASS because missing observations,
backend fallbacks and physical-audit warnings remain part of the
current evidence chain. Final submission should preserve these audit
records and should clearly state the definitive data source, station,
units, preprocessing decisions and calibration settings used to
reproduce the reported experiments.

%% =====================================================================
%% 6. CONCLUSIONS
%% =====================================================================
\section{Conclusions}\label{sec:conclusions}

This study evaluated StoHyMoLAP, an audit-controlled stochastic
physical--machine-learning framework for daily streamflow simulation in
the high-Andean Ramis basin. The central finding is that the main source
of robust improvement is not the stochastic perturbation itself or the
nominal complexity of the ML backend, but the use of quantile
descriptors from the stochastic physical ensemble as features for the
emulator.

The strongest backend-consistent hybrid, E6\_HYB\_XGB\_QUANTILES,
reached NSE\,=\,0.740 and KGE\,=\,0.810 over the common validation
period. This corresponds to a gain of 0.311 NSE relative to the full
stochastic physical baseline and 0.165 NSE relative to the pure
forcing-based ML baseline. Replacing ensemble-mean features with
ensemble quantiles alone added 0.114 NSE, confirming that the
distributional structure of the physical ensemble carries useful
information beyond a single mean trajectory.

The deterministic physical modifications had more limited effects. The
linear baseflow reservoir produced a modest improvement in global skill,
consistent with the need to represent slow catchment memory, but it did
not fully correct low-flow errors. Direct L\'{e}vy perturbations did not
reliably improve point-prediction skill when used only inside the
physical simulator. Their strongest contribution was indirect: they
created ensemble spread from which informative quantile descriptors
could be extracted for ML correction.

The uncertainty analysis shows that predictive skill and uncertainty
reliability must be assessed separately. The stochastic physical bands
were strongly underdispersive, with PICP\,=\,0.004 against a nominal
0.95 target. Thus, the present ensemble should not be interpreted as a
calibrated uncertainty forecast. It is currently better understood as a
feature-generation mechanism whose uncertainty representation requires
additional recalibration before operational use.

The backend audit changes the interpretation of the best numerical
score. E8\_HYB\_GRU\_QUANTILES achieved the highest NSE, but it was
executed as a fallback multilayer perceptron rather than as a true GRU.
Consequently, the defensible conclusion is not that recurrent networks
outperform the other models, but that quantile-based stochastic
physical descriptors consistently improve hybrid streamflow simulation
across the tested emulator families. Future work should rerun the
sequential experiments with verified recurrent backends, recalibrate the
stochastic ensemble for nominal coverage, incorporate low-flow-weighted
objectives and test the framework across additional Altiplano basins.

